% Finite-iteration root bracket with exact interval width.
\clearpage
\section*{Verification case: a root finder must certify its answer}
\noindent Find the positive real solution of \(f(x)=x^3-2=0\).
The function is continuous and strictly increasing on \([1,2]\);
\(f(1)=-1\) and \(f(2)=6\). Therefore at least one root lies in
the interval, and monotonicity makes it unique.

Bisection maintains an interval \([a_n,b_n]\) whose endpoints
bracket the root. After \(n\) successful halvings,
\[
b_n-a_n=\frac{b_0-a_0}{2^n}=\frac1{2^n}.
\]
Choosing its midpoint \(m_n\) gives the \emph{certified} bound
\[
|m_n-x_*|\le\frac1{2^{n+1}}.
\]
Ten halvings therefore guarantee a midpoint error no greater
than \(1/2048\). Numerically, this is approximately
\(0.00048828125\), independent of derivatives.
The exact root is \(2^{1/3}\approx1.25992105\).

Newton's method instead updates
\(x_{n+1}=x_n-(x_n^3-2)/(3x_n^2)\).
When started at \(x_0=0\), the derivative vanishes and this
formula is undefined. When started elsewhere, local rapid
convergence can be useful, but iteration count alone is not an
error certificate. Bracket safeguards can retain correctness while
allowing faster steps when their assumptions are satisfied.

\paragraph{Reconstruction.}
Implement bisection and a safeguarded Newton solver.
Return not only a floating-point estimate but also a final interval,
residual, iteration count and failure status.
Test exact endpoint roots, reversed endpoints, nonfinite values,
non-bracketing intervals, and a flat derivative at the initial guess.
Explain the difference between the width-based guarantee and a
residual that becomes small because \(f'\) is also small.
